\begin{graphillustration}{Граф $H$ из доказательства (здесь $n=6$, $k=2$): новые вершины $W$ образуют клику и смежны со всеми вершинами $G$. Совершенное паросочетание $H$ использует не более $k$ рёбер при $W$; остальные его рёбра — паросочетание $G$.}{fig:05-construction}
\begin{tikzpicture}[font=\small]
\node[draw,ellipse,minimum width=4.2cm,minimum height=2.2cm] (G) at (0,0) {};
\node[gi note] at (0,-1.4) {$G$};
\foreach \i/\x/\y in {1/-1.3/0.4,2/-0.3/0.6,3/0.9/0.5,4/1.2/-0.4,5/-1.1/-0.5,6/0.0/-0.5}
  \node[gi small] (g\i) at (\x,\y) {};
\node[gi vertex] (w1) at (4.2,0.8) {$w_1$};
\node[gi vertex] (w2) at (4.2,-0.8) {$w_2$};
\draw[gi edge] (w1) -- (w2);
\node[gi group,fit=(w1)(w2),label=right:{$W$, $|W|=k$}] {};
\foreach \i in {1,...,6} {\draw[gi edge,black!30] (w1) -- (g\i); \draw[gi edge,black!30] (w2) -- (g\i);}
\draw[gi match] (w1) -- (g3); \draw[gi match] (w2) -- (g4);
\draw[gi match] (g1) -- (g2); \draw[gi match] (g5) -- (g6);
\end{tikzpicture}
\end{graphillustration}
